\RequirePackage{plautopatch}
\documentclass[platex,dvipdfmx,a4paper]{jsarticle}

\usepackage{amsmath}
\usepackage{amsfonts}
\usepackage{amssymb}

\usepackage[dvipdfmx]{graphicx}
\usepackage{float}

\usepackage{tikz}
\usepackage{caption}

% 「%」は以降の内容を「改行コードも含めて」無視するコマンド
\usepackage[%
  dvipdfmx,% 欧文ではコメントアウトする
  setpagesize=false,%
  bookmarks=true,%
  bookmarksdepth=tocdepth,%
  bookmarksnumbered=true,%
  colorlinks=false,%
  pdftitle={},%
  pdfsubject={},%
  pdfauthor={},%
  pdfkeywords={}%
]{hyperref}
% PDFのしおり機能の日本語文字化けを防ぐ((u)pLaTeXのときのみかく)
\usepackage{pxjahyper}

\usepackage{framed}
\setlength{\OuterFrameSep}{0pt}

\usepackage{physics2}
\usephysicsmodule{ab}
\usephysicsmodule{ab.braket}


\newcommand{\hb}[2]{\frac{\partial #1}{\partial #2}}
\newcommand{\bbn}[2]{\frac{d #1}{d #2}}
\newcommand{\nhb}[3]{\frac{\partial^{#1} #2}{\partial #3^{#1}}}
\newcommand{\nbbn}[3]{\frac{d^{#1} #2}{d #3^{#1}}}
\newcommand{\vectxy}[2]{\begin{bmatrix}#1\\#2\end{bmatrix}}
\newcommand{\vectxyz}[3]{\begin{bmatrix}#1\\#2\\#3\end{bmatrix}}
\newcommand{\vect}[1]{\boldsymbol{\mathrm{#1}}}

\newcommand{\ssum}[3]{\sum_{#1=#2}^#3}

\newcommand{\veff}{V_{\mathrm{eff}}(r)}

\hypersetup{colorlinks=true}
\hypersetup{linkcolor=red}

\title{解析力学演習レポート2}
\author{25B00729　嶋田 悠音}

\begin{document}
\maketitle
\subsection*{(1)}

電流密度を

\begin{align*}
  \vect{J}(\vect{r},t)
  =
  I(t)\delta l\,\vect{e}_z\,\delta^{(3)}(\vect{r})
\end{align*}

とおく。遅延ベクトルポテンシャルは

\begin{align*}
  \vect{A}(\vect{r},t)
  =
  \frac{\mu_0}{4\pi}
  \int
  \frac{\vect{J}(\vect{r}',t-R/c)}{R}
  d^3r'
\end{align*}

で与えられる。ただし

\begin{align*}
  R=|\vect{r}-\vect{r}'|
\end{align*}

である。原点に局在した電流であるため、デルタ関数を用いて積分を実行すると

\begin{align*}
  \vect{A}(\vect{r},t)
  =
  \frac{\mu_0\delta l}{4\pi r}
  I\!\left(t-\frac{r}{c}\right)
  \vect{e}_z
\end{align*}

を得る。以後、

\begin{align*}
  t_r=t-\frac{r}{c}
\end{align*}

を遅延時刻とする。

ここで

\begin{align*}
  \vect{e}_z
  =
  \cos\theta\,\vect{e}_r
  -
  \sin\theta\,\vect{e}_\theta
\end{align*}

を用いると、

\begin{align*}
  A_r
  &=
  \frac{\mu_0\delta l}{4\pi r}
  I(t_r)\cos\theta,\\
  A_\theta
  &=
  -\frac{\mu_0\delta l}{4\pi r}
  I(t_r)\sin\theta
\end{align*}

となる。

ローレンツゲージ条件

\begin{align*}
  \nabla\cdot\vect{A}
  +
  \frac{1}{c^2}
  \hb{\phi}{t}
  =
  0
\end{align*}

を用いる。球座標における発散

\begin{align*}
  \nabla\cdot\vect{A}
  =
  \frac{1}{r^2}
  \hb{}{r}\!\left(r^2A_r\right)
  +
  \frac{1}{r\sin\theta}
  \hb{}{\theta}\!\left(\sin\theta A_\theta\right)
\end{align*}

に代入すると

\begin{align*}
  \nabla\cdot\vect{A}
  &=
  \frac{\mu_0\delta l\cos\theta}{4\pi}
  \left(
    -\frac{\dot I(t_r)}{cr}
    -\frac{I(t_r)}{r^2}
  \right)
\end{align*}

を得る。したがって

\begin{align*}
  \hb{\phi}{t}
  =
  \frac{\delta l\cos\theta}{4\pi\varepsilon_0}
  \left(
    \frac{\dot I(t_r)}{cr}
    +
    \frac{I(t_r)}{r^2}
  \right)
\end{align*}

であり、時間積分すると

\begin{align*}
  \phi
  =
  \frac{\delta l\cos\theta}{4\pi\varepsilon_0}
  \left(
    \frac{I(t_r)}{cr}
    +
    \frac{1}{r^2}
    \int^t I(t_r)\,dt
  \right)
\end{align*}

となる。

ここで電気双極子モーメント

\begin{align*}
  p(t)
  =
  \delta l
  \int^t I(t')\,dt'
\end{align*}

を導入すると

\begin{align*}
  \phi
  =
  \frac{\cos\theta}{4\pi\varepsilon_0}
  \left(
    \frac{\dot p(t_r)}{cr}
    +
    \frac{p(t_r)}{r^2}
  \right)
\end{align*}

と書ける。

磁場は

\begin{align*}
  \vect{B}
  =
  \nabla\times\vect{A}
\end{align*}

より求まる。軸対称性のため \(\phi\) 成分のみが残り、

\begin{align*}
  B_\phi
  &=
  \frac{1}{r}
  \left[
    \hb{}{r}(rA_\theta)
    -
    \hb{A_r}{\theta}
  \right] \\
  &=
  \frac{\mu_0\sin\theta}{4\pi}
  \left(
    \frac{\dot p(t_r)}{r^2}
    +
    \frac{\ddot p(t_r)}{cr}
  \right)
\end{align*}

を得る。

次に

\begin{align*}
  \vect{E}
  =
  -\nabla\phi
  -
  \hb{\vect{A}}{t}
\end{align*}

を計算する。

まず動径成分について、

\begin{align*}
  E_r
  &=
  -\hb{\phi}{r} \\
  &=
  \frac{\cos\theta}{2\pi\varepsilon_0}
  \left(
    \frac{p(t_r)}{r^3}
    +
    \frac{\dot p(t_r)}{cr^2}
  \right)
\end{align*}

を得る。

さらに極角成分は

\begin{align*}
  E_\theta
  &=
  -\frac{1}{r}\hb{\phi}{\theta}
  -\hb{A_\theta}{t} \\
  &=
  \frac{\sin\theta}{4\pi\varepsilon_0}
  \left(
    \frac{p(t_r)}{r^3}
    +
    \frac{\dot p(t_r)}{cr^2}
    +
    \frac{\ddot p(t_r)}{c^2r}
  \right)
\end{align*}

となる。

したがって求める電磁場は

\begin{align*}
  E_r
  &=
  \frac{\cos\theta}{2\pi\varepsilon_0}
  \left(
    \frac{p(t_r)}{r^3}
    +
    \frac{\dot p(t_r)}{cr^2}
  \right),\\
  E_\theta
  &=
  \frac{\sin\theta}{4\pi\varepsilon_0}
  \left(
    \frac{p(t_r)}{r^3}
    +
    \frac{\dot p(t_r)}{cr^2}
    +
    \frac{\ddot p(t_r)}{c^2r}
  \right),\\
  B_\phi
  &=
  \frac{\mu_0\sin\theta}{4\pi}
  \left(
    \frac{\dot p(t_r)}{r^2}
    +
    \frac{\ddot p(t_r)}{cr}
  \right)
\end{align*}

で与えられる。

\subsection*{(2)}

得られた電磁場を距離依存性で分類する。

まず

\begin{align*}
  \frac{p(t_r)}{r^3}
\end{align*}

に比例する項は静的双極子電場

\begin{align*}
  \vect{E}
  \propto
  \frac{p}{r^3}
\end{align*}

と同じ形をもつ。したがってこの項は電荷分布そのものに起因する静電場である。

次に

\begin{align*}
  \frac{\dot p(t_r)}{r^2}
\end{align*}

に比例する項は、時間変化する電流によって生じる誘導電磁場である。この項は静電場と放射場の中間的な性質をもち、源の近傍にエネルギーを蓄積したり戻したりする成分に対応する。

最後に

\begin{align*}
  \frac{\ddot p(t_r)}{r}
\end{align*}

に比例する項は放射電磁場である。実際、

\begin{align*}
  E_\theta
  =
  \frac{\sin\theta}{4\pi\varepsilon_0}
  \frac{\ddot p(t_r)}{c^2r},
  \qquad
  B_\phi
  =
  \frac{\mu_0\sin\theta}{4\pi}
  \frac{\ddot p(t_r)}{cr}
\end{align*}

より

\begin{align*}
  E_\theta
  =
  cB_\phi
\end{align*}

が成り立つ。また \(\vect{E}\)、\(\vect{B}\)、伝播方向 \(\vect{e}_r\) は互いに直交する。

さらにポインティングベクトル

\begin{align*}
  \vect{S}
  =
  \frac{1}{\mu_0}
  \vect{E}\times\vect{B}
\end{align*}

は

\begin{align*}
  S
  \propto
  \frac{1}{r^2}
\end{align*}

となる。したがって半径 \(r\) の球面を通過する全エネルギー流

\begin{align*}
  \int S\,dA
\end{align*}

は \(r\) に依らず一定であり、エネルギーが無限遠まで運ばれることが分かる。

以上より、\(1/r^3\) 項は静電場、\(1/r^2\) 項は誘導電磁場、\(1/r\) 項は電磁放射に対応することが示された。
\end{document}